\documentclass[10pt]{article}
\usepackage[T1]{fontenc}
\usepackage[utf8]{inputenc}
\usepackage[french]{babel}
\usepackage{amsmath,amssymb,mathrsfs,mathpazo}
\usepackage{geometry}
\geometry{a4paper,margin=2cm}
\usepackage{graphicx,array,tabularx,tabularray,colortbl,multirow,multicol,booktabs,enumitem,pifont,ulem,fancybox,diagbox,cancel}
\usepackage{pstricks,pstricks-add,pst-plot,pst-tree,pst-node,pst-text,pst-func,pst-math,pst-bspline,pst-3dplot,pst-eucl,pst-solides3d}
\let\pscalculate\undefined
\usepackage{pst-bezier}
\usepackage{tikz,pgfplots,tkz-tab}
\usetikzlibrary{arrows,arrows.meta,patterns,calc,positioning,trees,decorations.markings,shapes.geometric,shapes.misc,shapes.arrows,intersections,decorations.pathmorphing,decorations.pathreplacing,backgrounds,shadows,babel}
\pgfplotsset{compat=1.18}
\usepgfplotslibrary{fillbetween}
\usepackage[np]{numprint}
\usepackage{listings,makecell,pas-tableur,twoopt}
\let\barre\undefined
\usepackage[pstricks]{tablvar}
\definecolor{darkgreen}{RGB}{0,100,0}
\providecommand{\red}{\color{red}}
\providecommand{\blue}{\color{blue}}
\providecommand{\green}{\color{green}}
\providecommand{\black}{\color{black}}
\providecommand{\white}{\color{white}}
\providecommand{\gray}{\color{gray}}
\providecommand{\cyan}{\color{cyan}}
\providecommand{\magenta}{\color{magenta}}
\providecommand{\yellow}{\color{yellow}}
\providecommand{\violet}{\color{violet}}
\providecommand{\orange}{\color{orange}}
\providecommand{\euro}{\texteuro}
\providecommand{\EUR}{\texteuro}
\providecommand{\ieme}{\textsuperscript{e}}
\providecommand{\degres}{\ensuremath{^\circ}}
\providecommand{\degree}{\ensuremath{^\circ}}
\providecommand{\vv}[1]{\overrightarrow{#1}}
\providecommand{\Coord}[1]{#1}
\providecommand{\decosix}{}
\providecommand{\starredbullet}{\ensuremath{\bullet}}
\pagestyle{empty}
\setlength{\parindent}{0pt}
\usepackage{adjustbox}\DeclareUnicodeCharacter{2236}{:}\DeclareUnicodeCharacter{2190}{\ensuremath{\leftarrow}}
\providecommand{\R}{\mathbb{R}}
\providecommand{\N}{\mathbb{N}}
\providecommand{\Z}{\mathbb{Z}}
\providecommand{\Q}{\mathbb{Q}}
\providecommand{\D}{\mathbb{D}}
\providecommand{\C}{\mathbb{C}}
\providecommand{\e}{{\mathrm{e}}}
\providecommand{\vect}[1]{\overrightarrow{#1}}
\providecommand{\vectt}[1]{\overrightarrow{\mathrm{#1}}}
\providecommand{\vvt}[1]{\overrightarrow{\mathrm{#1}}}
\providecommand{\np}[2][]{{#2~#1}}
\providecommand{\vv}[1]{{\overrightarrow{#1}}}
\renewcommand{\d}{{\mathrm{d}}}
\providecommand{\degres}{{^\circ}}
\providecommand{\degre}{{^\circ}}
\providecommand{\blue}{\color{blue}}
\providecommand{\red}{\color{red}}
\providecommand{\enskip}{\quad}
\providecommand{\llbracket}{\lbrack\!\lbrack}
\providecommand{\rrbracket}{\rbrack\!\rbrack}
\providecommand{\euro}{{\text{€}}}
\providecommand{\no}{{\text{n°}}}
\providecommand{\interval}[2]{{[#1;#2]}}
\providecommand{\intervFO}[2]{{[#1;#2[}}
\providecommand{\dotfill}{\ldots}

\providecommand{\ds}{\displaystyle}
\providecommand{\textding}[1]{#1}
\providecommand{\Oij}{\(\mathrm{O};\vec{\imath},\vec{\jmath}\)}
\providecommand{\Oijk}{\(\mathrm{O};\vec{\imath},\vec{\jmath},\vec{k}\)}
\providecommand{\Ouv}{\(\mathrm{O};\vec{u},\vec{v}\)}
\providecommand{\intervFF}[2]{\ensuremath{[#1;#2]}}
\renewcommand{\euro}{\ifmmode\text{\texteuro}\else\texteuro\fi}
\DeclareUnicodeCharacter{00B0}{\ifmmode^{\circ}\else\textdegree\fi}
\DeclareUnicodeCharacter{20AC}{\ifmmode\text{\texteuro}\else\texteuro\fi}
% Énoncés : Éducation nationale. Transcription et corrigés : APMEP, Denis Vergès et contributeurs. Découpage et classement : CoopMaths, Sébastien Lozano, Éric Elter et contributeurs. Adaptation de la présentation : Maths Explorer.
% CC BY-SA : https://coopmaths.fr/www/about/
% https://coopmaths.fr/alea/static/eam/2026/tex/eam_technologique_2026_06_antilles_2.tex
\begin{document}
\noindent\textit{APMEP / CoopMaths. CC BY-SA. Présentation adaptée par Maths Explorer.}\par\medskip
En France, en 2021, une étude portant sur 100 personnes locataires du parc social (HLM) montre que :

	\begin{itemize}[label={\textbullet~}]
		\item 70 sont des personnes qui ne vivent pas en couple et parmi elles, la moitié est sans enfant ;
		\item Parmi les personnes vivant en couple, 12 n'ont pas d'enfant.
	\end{itemize}

	\begin{center}
		\begin{tblr}{|c|c|c|c|}
			\hline
			& sans enfant & au moins un enfant & Total \\
			\hline
			personne vivant en couple & 12 & 18 & 30 \\
			\hline
			personne ne vivant pas en couple & 35 & 35 & 70 \\
			\hline
			Total & 47 & 53 & 100 \\
			\hline
		\end{tblr}
	\end{center}

	On interroge une de ces personnes au hasard et on considère les évènements suivants :

	\begin{itemize}[label={\textbullet~}]
		\item $C$ : «\,la personne vit en couple\,»;
		\item $E$ : «\,la personne a au moins un enfant\,».
	\end{itemize}

	\begin{enumerate}
		\item Calculer la probabilité que la personne interrogée soit une personne sans enfant et ne vivant pas en couple.

		\item Calculer la probabilité de l'évènement $C \cap \overline{E}$ puis interpréter.

		\item Calculer la probabilité que la personne interrogée ne vive pas en couple sachant qu'elle a au moins un enfant.

		\item Dans cette question, on s'intéresse au nombre de pièces par logement dans le parc social.

		Une autre étude donne la répartition suivante :
		\begin{itemize}[label={\textbullet~}]
			\item $8\,\%$ des logements possèdent 5 pièces;
			\item $25\,\%$ ont 4 pièces;
			\item $40\,\%$ ont 3 pièces;
			\item $20\,\%$ ont deux pièces;
			\item et $7\,\%$ ont une seule pièce.
		\end{itemize}
		On choisit un logement au hasard et on note $X$ la variable aléatoire associant à chaque logement, son nombre de pièce(s).

		\begin{enumerate}
			\item Déterminer la loi de probabilité de $X$.

			\item Calculer $P(X \geqslant  2)$.

			Interpréter ce résultat dans le contexte de l'exercice.

			\item Montrer que $\mathrm{E}(X)=3,07$.

			Interpréter ce résultat.
		\end{enumerate}
	\end{enumerate}

\end{document}
